\section{Đánh giá kết quả thực hiện}

Phần này đánh giá mức độ hoàn thành so với yêu cầu mini project và giá trị của sản phẩm. Các mô tả chi tiết về từng API, màn hình, test case và kịch bản demo đã được trình bày ở các section trước nên không lặp lại tại đây.

\subsection{Mức độ đáp ứng yêu cầu}

Hệ thống đáp ứng nhóm yêu cầu bắt buộc: có báo cáo nghiên cứu, Backend API, Admin Dashboard, Demo Application, dữ liệu lưu trữ bền vững, validation, xử lý lỗi, seed data, tài liệu thiết kế và kịch bản trình diễn. Các chức năng cốt lõi của Feature Flag gồm quản lý flag, cấu hình rule, evaluation runtime và audit log đều hoạt động trong luồng end-to-end.

Ngoài MVP, hệ thống đã hoàn thành nhiều yêu cầu khuyến khích: cache evaluation snapshot, JavaScript SDK, unit/integration test cho rule engine và API, history dựa trên audit log, RBAC demo, statistics, group kill switch và Docker Compose.

\subsection{Đánh giá thiết kế kỹ thuật}

Điểm mạnh chính của thiết kế là tách rõ control plane và data plane. Control plane xử lý cấu hình, phân quyền và audit; data plane tập trung vào evaluation nhanh, xác định và fail-closed. Sự phân tách này giúp giải thích kiến trúc rõ ràng trong báo cáo và giảm rủi ro khi mở rộng runtime evaluation sau này.

Rule Evaluation Engine có precedence tường minh và reason code ổn định, giúp kết quả dễ kiểm thử. Mô hình dữ liệu tách định nghĩa flag khỏi cấu hình theo environment, hỗ trợ group kill switch và lưu audit log append-only. Dashboard và SDK bám theo contract của backend nên các thành phần phối hợp nhất quán.

\subsection{Đánh giá giá trị thực tế}

Sản phẩm chứng minh được giá trị thực tế của Feature Flag trong ba tình huống: phát hành tính năng theo phạm vi nhỏ trước khi mở rộng, rollback nhanh bằng kill switch và quan sát lịch sử thay đổi khi cấu hình runtime làm thay đổi hành vi ứng dụng. Demo app giúp người xem thấy trực tiếp sự khác biệt giữa deploy mã nguồn và release tính năng.

So với LaunchDarkly hoặc Unleash, mini project không hướng tới đầy đủ chức năng thương mại. Giá trị của hệ thống nằm ở việc tái hiện các nguyên lý cốt lõi bằng một kiến trúc nhỏ, dễ giải thích và có đủ bằng chứng triển khai.

\subsection{Hạn chế còn lại}

Hệ thống hiện vẫn là một bản mini project. Các hạn chế chính gồm: chưa đánh giá ở lưu lượng lớn, chưa có đồng bộ cấu hình realtime giữa nhiều instance, số loại rule còn giới hạn, SDK mới có JavaScript, mô hình xác thực chỉ phục vụ demo, và chưa có monitoring/backup/recovery ở mức production.

Những hạn chế này không làm thay đổi kết quả MVP nhưng xác định rõ các hướng phát triển cần ưu tiên nếu hệ thống tiếp tục được mở rộng sau chương trình.
